% t* 공간에서 분산의 미분 가능성
% 컴파일: xelatex 또는 tectonic (kotex, 맑은 고딕 사용)
\documentclass[11pt,a4paper]{article}

\usepackage[margin=25mm]{geometry}
\usepackage{amsmath,amssymb,amsthm}
\usepackage{kotex}
\setmainhangulfont{malgun.ttf}[Path=C:/Windows/Fonts/, BoldFont=malgunbd.ttf]
\setsanshangulfont{malgun.ttf}[Path=C:/Windows/Fonts/, BoldFont=malgunbd.ttf]
\usepackage{enumitem}
\usepackage{hyperref}

\newtheorem{theorem}{정리}
\newtheorem{lemma}{보조정리}
\theoremstyle{definition}
\newtheorem{definition}{정의}
\theoremstyle{remark}
\newtheorem{remark}{참고}

\newcommand{\Var}{\operatorname{Var}}
\newcommand{\R}{\mathbb{R}}

\title{$t^*$ 공간에서 분산의 미분 가능성}
\author{}
\date{2026년 11월 13일}

\begin{document}
\maketitle

\begin{abstract}
모든 내부 기저함수 $F_i^n$의 최댓값이 임계점 $t^*_i$ 하나에서만 나오는 노드들의 집합을 $T^\circ$라 하자.
$T^\circ$는 열린 집합이고, 그 위에서 최댓값들의 분산 $\Var(\tau)$는 노드 $\tau$의 유리함수, 즉 실해석적 함수임을 보인다.
최댓값이 두 후보에서 동시에 나오는 점(경계)은 이 결과에 포함되지 않는다.
\end{abstract}

\section{설정과 표기}

노드 $\tau=(\tau_1,\dots,\tau_n)$에 대해 다음 집합을 정의역으로 둔다.
\[
U=\bigl\{\tau\in\R^n:\ -1\le\tau_1<\tau_2<\dots<\tau_n\le 1,\ \ \omega_i^n(\tau)\neq0\ (i=1,\dots,n)\bigr\}.
\]
끝점을 $\tau_1=-1,\ \tau_n=1$로 고정하거나 대칭 매개변수 $(a,b,\dots)$를 쓰는 경우도, 매개변수에서 $\tau$로 가는 사상이 선형이므로 아래 결론이 그대로 성립한다.

\begin{itemize}[leftmargin=*]
\item 라그랑주 기저와 분모:
\[
L_i^n(t)=\prod_{j\ne i}\frac{t-\tau_j}{\tau_i-\tau_j},\qquad D_i^n=\prod_{j\ne i}(\tau_i-\tau_j).
\]
\item 가중치와 정규화된 적분:
\[
\Phi_i(t)=\int_{-1}^{t}L_i^n(s)\,ds,\qquad \omega_i^n=\Phi_i(1),\qquad
G_i(t)=\frac{\Phi_i(t)}{\omega_i^n}-\frac12 .
\]
따라서 $G_i(-1)=-\tfrac12$, $G_i(1)=\tfrac12$이다.
\item 내부 기저함수: $F_i^n(t)=G_i(t)-G_{i+1}(t)$ $(i=1,\dots,n-1)$. 모든 $i$에 대해 $F_i^n(\pm1)=0$이다.
\item 최댓값과 분산:
\[
M_i(\tau)=\max_{t\in[-1,1]}F_i^n(t;\tau),\qquad
\Var(\tau)=\frac1{n-1}\sum_{i=1}^{n-1}\bigl(M_i-\bar M\bigr)^2,\quad \bar M=\frac1{n-1}\sum_{i=1}^{n-1}M_i .
\]
\item 임계점 $t^*_i$: $\alpha_i=\omega_{i+1}^n/D_i^n,\ \beta_i=\omega_i^n/D_{i+1}^n$이고 $\alpha_i\ne\beta_i$일 때
\[
t^*_i=\frac{\alpha_i\tau_{i+1}-\beta_i\tau_i}{\alpha_i-\beta_i}.
\]
\end{itemize}

\section{보조정리}

\begin{lemma}[계수의 매끄러움]\label{lem:smooth}
$F_i^n(t;\tau)$는 $t$의 다항식이고, 그 계수는 $U$에서 $\tau$의 유리함수이다.
따라서 $(t,\tau)\mapsto F_i^n(t;\tau)$는 $\R\times U$에서 $C^\infty$이다.
\end{lemma}

\begin{proof}
노드가 서로 다르므로 $D_i^n\ne0$이고, $L_i^n$의 계수는 $\tau_j$들의 다항식을 $D_i^n$으로 나눈 것이다.
항별로 적분하면 $\Phi_i$의 계수와 $\omega_i^n=\Phi_i(1)$도 $\tau$의 유리함수이다.
$U$에서 $\omega_i^n\ne0$이므로 $G_i=\Phi_i/\omega_i^n-\tfrac12$의 계수도 유리함수이고, $F_i^n=G_i-G_{i+1}$도 마찬가지이다.
\end{proof}

\begin{lemma}[임계점]\label{lem:critical}
$F_i^n$의 임계점은 $\tau_j$ $(j\ne i,i+1)$와, $\alpha_i\ne\beta_i$일 때의 $t^*_i$뿐이다.
\end{lemma}

\begin{proof}
$\tilde L_i=L_i^n/\omega_i^n$이라 하면 $\frac{d}{dt}F_i^n=\tilde L_i-\tilde L_{i+1}$이다.
$P(t)=\prod_{j\ne i,i+1}(t-\tau_j)$로 두면
\[
\tilde L_i(t)-\tilde L_{i+1}(t)
=P(t)\left[\frac{t-\tau_{i+1}}{\omega_i^nD_i^n}-\frac{t-\tau_i}{\omega_{i+1}^nD_{i+1}^n}\right]
=\frac{P(t)}{\omega_i^n\omega_{i+1}^n}\Bigl[\alpha_i(t-\tau_{i+1})-\beta_i(t-\tau_i)\Bigr].
\]
대괄호 안은 $(\alpha_i-\beta_i)t-(\alpha_i\tau_{i+1}-\beta_i\tau_i)$이다.
$\alpha_i\ne\beta_i$이면 근은 $t^*_i$ 하나이다.
$\alpha_i=\beta_i$이면 대괄호는 0이 아닌 상수 $\alpha_i(\tau_i-\tau_{i+1})$이므로 근이 없다.
\end{proof}

\begin{lemma}[후보 집합]\label{lem:candidates}
$C_i(\tau)=\{\pm1\}\cup\{\tau_j: j\ne i,i+1\}\cup\{t^*_i:\ \alpha_i\ne\beta_i,\ t^*_i\in[-1,1]\}$라 하면
\[
M_i(\tau)=\max_{c\in C_i(\tau)}F_i^n(c;\tau).
\]
\end{lemma}

\begin{proof}
$F_i^n$은 닫힌 구간 $[-1,1]$에서 연속이므로 최댓값은 끝점이나 내부 임계점에서 나온다.
보조정리~\ref{lem:critical}에 의해 내부 임계점은 모두 $C_i$에 들어 있다.
\end{proof}

\begin{lemma}[양수 가중치에서 $t^*_i$]\label{lem:positive}
모든 $\omega_j^n>0$이면 $\alpha_i\ne\beta_i$이고 $t^*_i\in(\tau_i,\tau_{i+1})$이다.
\end{lemma}

\begin{proof}
$D_i^n$에는 음수 인자 $\tau_i-\tau_j$ $(j>i)$가 $n-i$개 있으므로 $\operatorname{sgn}D_i^n=(-1)^{n-i}$이고, $\operatorname{sgn}D_{i+1}^n=(-1)^{n-i-1}$이다.
따라서 $\alpha_i$와 $\beta_i$는 부호가 반대이고 $\alpha_i-\beta_i\ne0$이다.
\[
t^*_i=\frac{\alpha_i}{\alpha_i-\beta_i}\,\tau_{i+1}+\frac{-\beta_i}{\alpha_i-\beta_i}\,\tau_i
\]
에서 두 계수는 모두 양수이고 합이 1이므로, $t^*_i$는 $\tau_i$와 $\tau_{i+1}$의 내분점이다.
\end{proof}

\begin{lemma}[유한 개 함수의 최댓값]\label{lem:max}
$f_1,\dots,f_r$가 열린 집합 $V$에서 연속이고, 점 $p\in V$에서 $f_1(p)>f_j(p)$ $(j\ne1)$라 하자.
그러면 $p$의 어떤 근방에서 $\max_j f_j=f_1$이다.
\end{lemma}

\begin{proof}
$\delta=\min_{j\ne1}\bigl(f_1(p)-f_j(p)\bigr)>0$이다.
연속성에 의해 $p$의 근방 $W$에서 모든 $j$에 대해 $|f_j-f_j(p)|<\delta/2$이다.
그러면 $W$에서 $f_1>f_1(p)-\delta/2\ge f_j(p)+\delta/2>f_j$이다.
\end{proof}

\section{정리}

\begin{definition}[$t^*$ 공간의 내부]
$g_i(\tau)=F_i^n\bigl(t^*_i(\tau);\tau\bigr)$라 하자. 다음을 만족하는 $\tau\in U$의 집합을 $T^\circ$라 한다.
모든 $i=1,\dots,n-1$에 대해
\begin{enumerate}[label=(\roman*)]
\item $\alpha_i(\tau)\ne\beta_i(\tau)$이고 $t^*_i(\tau)\in[-1,1]$이다.
\item $t^*_i$가 최댓값을 주는 \emph{유일한} 후보이다. 즉 $g_i(\tau)>F_i^n(c;\tau)$ $\ \forall c\in C_i(\tau)\setminus\{t^*_i\}$.
\end{enumerate}
\end{definition}

\begin{theorem}\label{thm:main}
$T^\circ$는 열린 집합이고, $T^\circ$에서 $M_i=g_i$ $(i=1,\dots,n-1)$이다.
따라서 $\Var$는 $T^\circ$에서 $\tau$의 유리함수이고, 특히 $C^\infty$이며 실해석적이다.
\end{theorem}

\begin{proof}
\textbf{(1) $g_i$의 매끄러움.}
$\alpha_i\ne\beta_i$인 곳은 열린 집합이고, 그곳에서 $t^*_i$는 $\tau$의 유리함수이다.
보조정리~\ref{lem:smooth}에 의해 합성 $g_i(\tau)=F_i^n(t^*_i(\tau);\tau)$도 유리함수이다.
후보 $\tau_j$에서의 값 $F_i^n(\tau_j;\tau)$와 $F_i^n(\pm1)=0$도 같은 이유로 유리함수이다.

\textbf{(2) 후보 집합이 국소적으로 바뀌지 않는다.}
$p\in T^\circ$라 하자. 조건 (ii)에서 $c=\pm1$을 고르면 $g_i(p)>F_i^n(\pm1)=0$이므로 $t^*_i(p)\neq\pm1$, 즉 $t^*_i(p)\in(-1,1)$이다.
$t^*_i$는 연속이므로 $p$의 근방에서 $\alpha_i\ne\beta_i$이고 $t^*_i\in(-1,1)$이다.
따라서 이 근방에서는 후보 집합 $C_i$에 같은 원소들(같은 이름의 후보)이 계속 들어 있다.

\textbf{(3) 열린 집합.}
(2)의 근방에서 조건 (ii)는 연속함수에 대한 유한 개의 강한 부등식이므로, 보조정리~\ref{lem:max}에 의해 더 작은 근방에서도 성립한다.
조건 (i)도 (2)에서 보였다. 따라서 $p$의 근방 전체가 $T^\circ$에 속하고, $T^\circ$는 열린 집합이다.

\textbf{(4) $M_i=g_i$.}
보조정리~\ref{lem:candidates}와 보조정리~\ref{lem:max}에 의해 $p$의 근방에서 $M_i=\max_{c\in C_i}F_i^n(c;\cdot)=g_i$이다.

\textbf{(5) 분산.}
$\Var$는 $M_1,\dots,M_{n-1}$의 다항식이므로 $T^\circ$에서 $g_i$들의 다항식, 즉 $\tau$의 유리함수이다.
유리함수는 분모가 0이 아닌 곳에서 실해석적이다.
\end{proof}

\begin{remark}[양수 가중치는 필요하지 않다]
정리~\ref{thm:main}의 증명은 가중치의 부호를 쓰지 않는다.
양수 가중치는 보조정리~\ref{lem:positive}에 의해 조건 (i)을 자동으로 보장해 줄 뿐이다.
가중치에 음수가 있어도 $t^*_i$가 유일한 최대이기만 하면 그 근방에서 $\Var$는 매끄럽다.
\end{remark}

\begin{remark}[편미분]
$T^\circ$에서 $t^*_i$는 $F_i^n$의 내부 임계점이므로, 포락선 정리에 의해
$\partial M_i/\partial\tau_k=(\partial F_i^n/\partial\tau_k)(t^*_i)$이다.
즉 $t^*_i$가 $\tau$에 따라 어떻게 움직이는지는 계산할 필요가 없다.
\end{remark}

\section{증명하지 않은 것}

\begin{enumerate}[leftmargin=*]
\item \textbf{경계(동점).} $t^*_i$와 다른 후보가 같은 최댓값을 주는 점은 $T^\circ$에 속하지 않는다.
두 매끄러운 함수 $g,h$에 대해 $\max(g,h)$는 $g=h$인 점에서 $\nabla g=\nabla h$가 아니면 미분 불가능하다.
따라서 동점을 포함한 닫힌 ``$t^*$ 공간'' 전체에서 미분 가능하다는 주장은 성립하지 않는다.
\item \textbf{$T^\circ$의 모양.} $T^\circ$가 연결인지, 양수 가중치 영역에 포함되는지는 이 정리로 알 수 없다.
\item \textbf{수치 검증.} 편미분 공식과 정확한 유한차분이 일치한 것은 계산 실수가 없다는 증거일 뿐, 증명은 아니다.
\end{enumerate}

\section{Newton 방법과의 관계}

Newton 방법이 필요로 하는 것은 해 근처에서의 매끄러움뿐이며, 정리~\ref{thm:main}이 이를 보장한다.
해가 $T^\circ$에 들어 있는지는 계산이 끝난 뒤 두 가지로 확인할 수 있다.
모든 가중치가 양수인지, 그리고 모든 $i$에서 $t^*_i$가 엄격하게 최대인지 보면 된다.
$n=4,\dots,21$에서 ESC 노드로 시작한 수렴점은 모두 이 조건을 만족했다.
$t^*_i$에서의 값과 두 번째로 큰 후보 값의 차이는 약 $0.76$으로, 수렴점이 $T^\circ$의 경계에서 충분히 떨어져 있었다.

\end{document}
